% displayname :
% !TEX root = /home/eb/Dropbox/CPGE/Physique/Cours/[5]Electromg/[2]ExempleElec/Doc/Docvdw.tex
\input{/home/eb/Dropbox/texmf/tex/latex/Preambles/TD.sty}
\begin{document}
\EnteteDoc{5}{2}

%%%%%%%%%%%%%%%%%%%%%%%%%%%%%%%%%%%%%%%%%%%%%%%%%%%%%%%%%%%%%%%%%%%%%%%%%%%%%%%%%%%%%%%%%%


%%% EXERCICE %%%
\begin{doc}{Les forces de Van der Waals et le Gecko}
%\begin{quest}
\item Quelles sont les trois types d'interactions de type Van der Waals ? De quelles natures sont ces interactions ?

\sol{
Un cadre de référence est un référentiel. Voir cours.
}

\item Pourquoi l'interaction de Keesom dépend-elle de la température ?

\sol{
Un référentiel inertiel est un référentiel dans lequel la 2e loi le Newton est valide.
}

\item À partir des ordres de grandeur d'énergie donnés, et en estimant la distance typique entre molécules, donner l'ordre de grandeur des forces de Van der Waals.

\sol{
Il s'agit de la force d'inertie d'entraînement, qui est axifuge lors d'une rotation uniforme.
Le ballon se déplace dans le référentiel du lanceur sur le carrousel dans une direction non colinéaire au vecteur rotation, et subit donc aussi la force de Coriolis.
}

\item Dans la partie 2.3, pourquoi passe-t-on d'une force en $1/r^7$ en une force en $1/r^3$ lorsque deux plans s'attirent ?

\sol{
On utilise la règle de la main droite.
%Dans la phrase de la p.3, l'auteur parle de la droite et de la gauche par rapport à un observateur qui observe le mouvement sur une carte.
}

\item Expliquer et vérifier les calculs de la partie 2.3.
\sol{
La force d'inertie d'entraînement contribue au poids, qui est vertical et non n'apparaît pas dans une projection horizontale. Avec $\Vec{\Omega}_{\rm T} = \Omega(\cos \phi \ey + \sin \phi \ez)$, la force de Coriolis s'écrit
$$
\Vec{f}_{\rm C} = 2\rho \Vec{v} \pv \Vec{\Omega}_{\rm T} = 2 \rho \Omega  (v_x \ex + v_y \ey) \pv (\cos \phi \ey + \sin \phi \ez) = \frac{4\pi\rho}{T}(v_x \cos \phi\, \ez - v_x \sin \phi\, \ey + v_y \sin \phi\,\ex) ,
$$
soit $\Vec{f}_{\rm C}  = \rho f(-v_x \ey + v_y \ex) + \text{composante verticale},$
avec $f = \frac{4\pi \sin \phi}{T}$.
}

\end{doc}
\end{document}
